% Chapter 17 -- Four Objects Are One (key: pseudorandom)
\newcommand{\pseudorandomList}{\mathrm{LIST}}
\newcommand{\pseudorandomAgr}{\operatorname{agr}}
\newcommand{\pseudorandomCP}{\operatorname{CP}}
\newcommand{\pseudorandomKL}{\operatorname{KL}}
\newcommand{\pseudorandomExt}{\operatorname{Ext}}
\newcommand{\pseudorandomEnc}{\operatorname{Enc}}
\newcommand{\pseudorandomSamp}{\operatorname{Samp}}
\newcommand{\pseudorandomHmin}{H_{\infty}}
\newcommand{\pseudorandomZigzag}{\mathbin{\text{\textcircled{\scriptsize z}}}}

\casechapter{Four Objects Are One}{Expanders, extractors, codes, and samplers}{Unification}{ch:pseudorandom}
\begin{refsection}[chapters/17-pseudorandom.bib]

\begin{chaptersummary}
Four communities built four kinds of pseudorandom object (expander graphs, list-decodable codes, randomness extractors, and samplers), each with its own definitions, existence proofs, and constructions. The move is to write all of them as one bipartite graph $\Gamma\colon[N]\times[D]\to[M]$ and every property as a bound on the size of one set: the left vertices that a test fools. Then one Cauchy--Schwarz inequality proves the expander mixing lemma, the Johnson bound, and the leftover hash lemma, one union bound produces the random objects, and constructions travel between the fields: a pseudorandom generator becomes an extractor, extractors become optimal samplers, and an algebraic code becomes an expander.
\end{chaptersummary}

\section{Before: four objects, four literatures}\label{sec:pseudorandom:before}

Randomized algorithms consume random bits, which are neither free nor perfect. Four communities asked four questions about this resource, and we meet their answers in their native languages.

\subsection{Expanders}\label{sec:pseudorandom:expanders}

A bipartite multigraph with left vertices $[N]=\{1,\dots,N\}$, right vertices $[M]$, and left degree $D$ is given by its \term{neighbor function} $\Gamma\colon[N]\times[D]\to[M]$: the value $\Gamma(x,y)$ is the $y$th neighbor of $x$, with repetitions allowed. For $S\subseteq[N]$ let $\Gamma(S)$ be the set of neighbors of vertices in $S$.

\begin{definition}[Vertex expander]\label{def:pseudorandom:expander}
Let $K\ge1$ be an integer and $A>0$. The graph $\Gamma$ is a \emph{$(K,A)$ vertex expander} if $|\Gamma(S)|\ge A|S|$ for every $S$ with $|S|\le K$, and a \emph{$({=}K,A)$ vertex expander} if $|\Gamma(S)|\ge AK$ whenever $|S|=K$.
\end{definition}

Expanders\index{expander graph} entered mathematics through telephone switching.\sidenote{Pinsker's paper appeared at a teletraffic conference, a meeting on the engineering of telephone networks.} \textcite{Pinsker1973} proved that random graphs expand; in the form of \textcite[Theorem~4.4]{Vadhan2012}, for every $D\ge3$ there is $\alpha>0$ such that a random $\Gamma\colon[N]\times[D]\to[N]$, with independent uniform values, is an $(\alpha N,D-2)$ vertex expander with probability at least $1/2$, by a count of repeated values among the neighbors of a $K$-set. Explicit constructions, starting with \textcite{Margulis1973}, needed deeper mathematics, and the most useful measure of expansion turned out to be spectral.\index{expander graph!spectral expansion} For a $D$-regular undirected multigraph on $[N]$ the random-walk matrix $P$, with $P_{xz}=\Pr_y[\Gamma(x,y)=z]$ for uniform $y\in[D]$, is symmetric and fixes the all-ones vector $\mathbf1$; let $\lambda(\Gamma)$ be the largest value of $\norm{Pf}/\norm f$ over nonzero $f\perp\mathbf1$. The \term{Ramanujan graphs}\index{Ramanujan graph} of \textcite{LubotzkyPhillipsSarnak1988} include families with $\lambda\le2\sqrt{D-1}/D$, asymptotically the least possible value.

\begin{theorem}[Expander mixing lemma]\label{thm:pseudorandom:mixing}
Let $\Gamma$ be a $D$-regular undirected multigraph on $[N]$ with $\lambda=\lambda(\Gamma)$. For $S,T\subseteq[N]$ let $e(S,T)$ be the number of pairs $(x,y)\in S\times[D]$ with $\Gamma(x,y)\in T$, and let $\sigma=|S|/N$ and $\tau=|T|/N$. Then
\[
  \Bigl|\frac{e(S,T)}{ND}-\sigma\tau\Bigr|\le\lambda\sqrt{\sigma(1-\sigma)\tau(1-\tau)}.
\]
\end{theorem}

Forms of the lemma\index{expander mixing lemma} go back to \textcite{AlonChung1988}; the textbook proof expands $\ind_S$ and $\ind_T$ in eigenvectors of $P$.

\subsection{List-decodable codes}\label{sec:pseudorandom:codes}

A \term{code} $\pseudorandomEnc\colon[N]\to[q]^D$ sends messages to words over the alphabet $[q]$; the \term{agreement} $\pseudorandomAgr(u,v)$ of two words is the fraction of coordinates where they coincide. Beyond half the minimum distance a decoder can only return a list.

\begin{definition}[List decoding]\label{def:pseudorandom:listdec}
Let $\rho\in(0,1)$ and let $L\ge1$ be an integer. The code is \emph{$(\rho,L)$ list-decodable} if for every $r\in[q]^D$ at most $L$ messages $x$ satisfy $\pseudorandomAgr(\pseudorandomEnc(x),r)>1-\rho$.
\end{definition}

List decoding\index{list decoding} was proposed by Elias and, independently, by Wozencraft in the late 1950s \parencite[\S5.6]{Vadhan2012}. With the $q$-ary entropy $H_q(\rho)\defeq\rho\log_q(q-1)-\rho\log_q\rho-(1-\rho)\log_q(1-\rho)$, random codes reach the following rate \parencite{Elias1991}, optimal up to the dependence on $L$ \parencite[\S5.1]{Vadhan2012}.

\begin{theorem}[Random codes are list-decodable]\label{thm:pseudorandom:elias}
Let $q\ge2$ and $L\ge1$ be integers, $\rho\in(0,1-1/q)$, and $R<1-H_q(\rho)-1/(L+1)$. For every $D$ and every $N\le q^{RD}$, a random code $\pseudorandomEnc\colon[N]\to[q]^D$ with independent uniform codewords is $(\rho,L)$ list-decodable with positive probability.
\end{theorem}

The classical proof is a union bound with the volume of a Hamming ball as its key estimate. For explicit codes, the distance bounds the lists.

\begin{theorem}[Johnson bound]\label{thm:pseudorandom:johnson}
Let $\eta\ge0$ and suppose that $\pseudorandomAgr(\pseudorandomEnc(x),\pseudorandomEnc(x'))\le1/q+\eta$ for all messages $x\ne x'$. If $\eps>\sqrt{(1-1/q)\eta}$, then for every $r\in[q]^D$ fewer than $(1-1/q)^2/(\eps^2-(1-1/q)\eta)$ messages $x$ satisfy $\pseudorandomAgr(\pseudorandomEnc(x),r)>1/q+\eps$.
\end{theorem}

Bounds of this kind\index{Johnson bound} go back to \textcite{Johnson1962}. Algorithms followed: \textcite{Sudan1997} list-decoded Reed--Solomon codes beyond half their distance, and \textcite{GuruswamiSudan1999} reached a $1-\sqrt R$ fraction of errors at rate $R$. None of this mentioned graphs or random sources.

\subsection{Weak sources and extractors}\label{sec:pseudorandom:extractors}

A random variable $X$ on $[N]$ has \term{min-entropy}\index{min-entropy} $\pseudorandomHmin(X)\defeq\min_x\log(1/\Pr[X=x])$, with logarithms to base $2$; it is a \term{$k$-source} if no value has probability above $2^{-k}$, and a \term{flat} $k$-source if it is uniform on $2^k$ values. The \term{statistical distance}\index{statistical distance} of random variables $Z$ and $W$ on $[M]$, $\Delta(Z,W)\defeq\max_{T\subseteq[M]}|\Pr[Z\in T]-\Pr[W\in T]|$, equals $\frac12\sum_z|\Pr[Z=z]-\Pr[W=z]|$, and $Z$ is $\eps$-close to $W$ if $\Delta(Z,W)\le\eps$. Santha and Vazirani showed in 1986 that no deterministic procedure extracts a nearly unbiased bit from every source in a natural class of weak sources \parencite[\S6.5]{Vadhan2012}; the way out, formalized by \textcite{NisanZuckerman1996}, is to invest a few truly random bits.

\begin{definition}[Extractor]\label{def:pseudorandom:extractor}
Let $Y$ be uniform on $\{0,1\}^d$. A function $\pseudorandomExt\colon\{0,1\}^n\times\{0,1\}^d\to\{0,1\}^m$ is a \emph{$(k,\eps)$ extractor} if $\pseudorandomExt(X,Y)$ is $\eps$-close to uniform for every $k$-source $X$ independent of $Y$, and \emph{strong} if moreover $(Y,\pseudorandomExt(X,Y))$ is $\eps$-close to uniform. The variable $Y$ is the \emph{seed}.
\end{definition}

The first extractors\index{randomness extractor} came from hashing. A family $\cH$ of functions $h\colon[N]\to[M]$ is \term{universal} if $\Pr_h[h(x)=h(x')]\le1/M$ for all $x\ne x'$, where $h$ is uniform on $\cH$. The \term{collision probability}\index{collision probability} $\pseudorandomCP(Z)\defeq\sum_z\Pr[Z=z]^2$ is the probability that two independent copies of $Z$ coincide.

\begin{theorem}[Leftover hash lemma]\label{thm:pseudorandom:lhl}
Let $\cH$ be a universal family of functions $[N]\to[M]$, let $X$ be a $k$-source on $[N]$, and let $\eps\in(0,1)$ satisfy $M\le\eps^22^k$. If $h$ is uniform on $\cH$ and independent of $X$, then $(h,h(X))$ is $\eps/2$-close to uniform on $\cH\times[M]$.
\end{theorem}

Forms of the lemma\index{leftover hash lemma} are due to Bennett, Brassard, and Robert in 1988 and to \textcite{ImpagliazzoLevinLuby1989}; the following proof is credited to Rackoff \parencite[\S6.5]{Vadhan2012}.\sidenote{In cryptography the lemma turns a partly secret string into a nearly uniform key; see \cref{ex:pseudorandom:privacy}.}

\begin{proof}
The plan is to bound the collision probability, convert it into an $\ell_2$ distance, and convert that into a statistical distance. Let $Z=(h,h(X))$ and let $Z'=(h',h'(X'))$ be an independent copy. They coincide only if $h=h'$ and either $X=X'$, which has probability $\pseudorandomCP(X)\le2^{-k}$, or $X\ne X'$ and the hash values collide, which has probability at most $1/M$ by universality; hence $\pseudorandomCP(Z)\le(2^{-k}+1/M)/|\cH|$. Identify $Z$ with its distribution vector, let $W=|\cH|M$, and let $u$ be the uniform distribution on $\cH\times[M]$; then $\norm{Z-u}_2^2=\pseudorandomCP(Z)-1/W\le2^{-k}/|\cH|$. By the Cauchy--Schwarz inequality over the $W$ coordinates, $\Delta(Z,u)=\frac12\norm{Z-u}_1$ is at most $\frac12\sqrt W\,\norm{Z-u}_2\le\frac12\sqrt{M2^{-k}}\le\eps/2$.
\end{proof}

The seed describes $h$, which takes $\Omega(n)$ bits, where $n=\log N$ \parencite[\S6.2]{Vadhan2012}. A simulation enumerates all seeds, so one wants seeds of $O(\log n)$ bits, and getting there took a decade.

\begin{aside}[Extractors before 1999]
The constructions of the 1990s worked in stages. \emph{Block sources}: if every block of $X=(X_1,\dots,X_t)$ keeps min-entropy given the earlier ones, extract from the last block with a short seed and use the output as the seed for the block before it. \emph{Reduction to block sources}: restrict $X$ to a pseudorandom set of bit positions that keeps a constant fraction of the entropy rate, paying $\log(1/\eps)$ bits each time through a chain rule for min-entropy. \emph{Recursion}, with errors and losses tracked level by level, handles lower entropy rates. \textcite{NisanZuckerman1996} obtained the first extractors with polylogarithmic seeds this way, and later work lowered the entropy with more elaborate recursions \parencite[\S6.3 and \S6.5]{Vadhan2012}. Each stage has its own lemma, and nothing suggests that an extractor and a list-decodable code are the same kind of object.
\end{aside}

\subsection{Samplers and amplification}\label{sec:pseudorandom:samplers}

The fourth question is estimation: given oracle access to $f\colon\{0,1\}^m\to[0,1]$, estimate its mean within $\eps$, with failure probability at most $\delta$, using few random bits. Bellare and Rompel isolated the form in which the queries depend only on the random bits and the estimate is their average.\sidenote{They, and \textcite{Zuckerman1997}, called these oblivious samplers \parencite{Vadhan2012}.} Let $\mu(f)$ be the average of $f\colon[M]\to[0,1]$, and let $x$ be uniform on $[N]$.

\begin{definition}[Averaging sampler]\label{def:pseudorandom:sampler}
A function $\pseudorandomSamp\colon[N]\to[M]^D$ is a \emph{$(\delta,\eps)$ averaging sampler} if for every $f\colon[M]\to[0,1]$
\[
  \Pr_x\Bigl[\frac1D\sum_{y=1}^Df\bigl(\pseudorandomSamp(x)_y\bigr)>\mu(f)+\eps\Bigr]\le\delta ,
\]
and a \emph{Boolean} one if this holds for indicator functions. It uses $\log N$ random bits and makes $D$ queries.
\end{definition}

Applying the definition to $1-f$ bounds the other tail, at the price of a factor $2$ in $\delta$. Independent samples need $O(\log(1/\delta)/\eps^2)$ queries\index{averaging sampler} but $m$ random bits each; pairwise-independent samples $ai+b$, for $D$ values of $i$ and uniform $(a,b)\in\F_{2^m}^2$, need $2m$ bits, but Chebyshev's inequality then gives only $\delta=1/(4D\eps^2)$, as Chor and Goldreich observed \parencite[\S3.5]{Vadhan2012}. Reducing the error of a randomized algorithm is the same problem, since a majority vote errs only if the runs over-represent the bad random strings. Walks on expanders, analyzed in a line of work starting with Ajtai, Koml\'os, and Szemer\'edi in 1987, reduce the error to $2^{-t}$ with $m+O(t)$ random bits, through a spectral analysis of constrained walks or the Chernoff bound for walks of \textcite{Gillman1998} \parencite[\S4.6]{Vadhan2012}. \textcite{Sipser1988} built an early bridge between two of our worlds: samplers that hit every large set are vertex expanders.

Some coincidences were hard to miss. Random samplers, extractors, and binary codes at radius $1/2-\eps$ all need about $\log(\cdot)/\eps^2$ queries, seeds, or coordinates, and the mixing lemma, the Johnson bound, and the leftover hash lemma are all a second moment followed by Cauchy--Schwarz.

\section{The move: one graph, one list}\label{sec:pseudorandom:move}

Write every object as a function $\Gamma\colon[N]\times[D]\to[M]$ of two arguments. The first, $x$, is the randomness we pay for: the coins of a sampler, a sample from a weak source, a message, a left vertex. The second, $y$, is enumerated: a query index, a seed, a coordinate, an edge label.\sidenote{For this reason \textcite[\S6.2]{Vadhan2012} names the seed length $d$: an extractor is a graph of left degree $D=2^d$.} The value $\Gamma(x,y)$ is a sample, an output, a coordinate with its symbol, a neighbor. Each property of \cref{sec:pseudorandom:before} then says that no test is fooled by many values of $x$.

\begin{definition}[Lists]\label{def:pseudorandom:list}
Let $\Gamma\colon[N]\times[D]\to[M]$. For $x\in[N]$ and $T\subseteq[M]$ let $\Gamma_x(T)\defeq\Pr_y[\Gamma(x,y)\in T]$, with $y$ uniform on $[D]$, and let $\mu(T)\defeq|T|/M$. For $\theta\in[0,1)$ the \emph{list} of the test $T$ at threshold $\theta$ is
\[
  \pseudorandomList_\Gamma(T,\theta)\defeq\{x\in[N]:\Gamma_x(T)>\theta\},
\]
and $\pseudorandomList_\Gamma(T,1)\defeq\{x:\Gamma_x(T)=1\}$. For $f\colon[M]\to[0,1]$ define $\Gamma_x(f)\defeq\E_yf(\Gamma(x,y))$ and $\pseudorandomList_\Gamma(f,\theta)$ in the same way. The list at a threshold $\mu(T)+\eps\ge1$ or $\mu(f)+\eps\ge1$ is empty.
\end{definition}

The new object is humble: a set of bad inputs, counted.\index{list of a test} This is Grothendieck's advice to stoop low: pick up the count that every definition already contained. The list forgets which inputs are fooled, and how; it keeps only how many, which is all that each of the four definitions uses. Two facts about tests come first.

\begin{lemma}[Tests]\label{lem:pseudorandom:tests}
Let $Z$ and $W$ be random variables on $[M]$.
\begin{enumerate}
\item The distance satisfies $\Delta(Z,W)=\max_T\bigl(\Pr[Z\in T]-\Pr[W\in T]\bigr)$, and every $f\colon[M]\to[0,1]$ satisfies $|\E f(Z)-\E f(W)|\le\Delta(Z,W)$.
\item Let $K=2^k$ be an integer. A function $\Gamma\colon[N]\times[D]\to[M]$ is a $(k,\eps)$ extractor if and only if $\Pr_{x\in S,\,y}[\Gamma(x,y)\in T]\le\mu(T)+\eps$ for every $S\subseteq[N]$ with $|S|=K$ and every $T\subseteq[M]$, where $x$ is uniform on $S$.
\end{enumerate}
\end{lemma}

\begin{proof}
(i) Replacing $T$ by its complement changes the sign of the difference, so the largest signed difference is the largest absolute one. For $f$ with values in $[0,1]$, the difference $\E f(Z)-\E f(W)$ equals $\int_0^1(\Pr[f(Z)>s]-\Pr[f(W)>s])\dd s\le\Delta(Z,W)$, and symmetrically.

(ii) The condition is necessary because the uniform distribution on a set of size $K$ is a $k$-source. For sufficiency, let $X$ be a $k$-source and $T\subseteq[M]$. The probability $\Pr[\Gamma(X,y)\in T]=\sum_x\Pr[X=x]\,\Gamma_x(T)$ averages the numbers $\Gamma_x(T)$ with weights in $[0,1/K]$. Moving weight toward larger numbers does not decrease it, so it is at most the average of the $K$ largest numbers $\Gamma_x(T)$, which is $\Pr_{x\in S,y}[\Gamma(x,y)\in T]$ for the set $S$ of the corresponding $x$. Part (i) turns the bound for every $T$ into $\Delta\le\eps$.
\end{proof}

\begin{proposition}[The dictionary]\label{prop:pseudorandom:dictionary}
Let $\Gamma\colon[N]\times[D]\to[M]$.
\begin{enumerate}
\item \emph{Expanders.} For an integer $K\ge1$ and $A>0$, $\Gamma$ is a $({=}K,A)$ vertex expander if and only if $|\pseudorandomList_\Gamma(T,1)|<K$ for every $T\subseteq[M]$ with $|T|<AK$.
\item \emph{Codes.} Let $M=Dq$, identify $[M]$ with $[D]\times[q]$, and let $\Gamma(x,y)=(y,\pseudorandomEnc(x)_y)$. For $r\in[q]^D$ let $T_r\defeq\{(y,r_y):y\in[D]\}$. Then $\Gamma_x(T_r)=\pseudorandomAgr(\pseudorandomEnc(x),r)$, so $\pseudorandomEnc$ is $(\rho,L)$ list-decodable if and only if $|\pseudorandomList_\Gamma(T_r,1-\rho)|\le L$ for every $r$.
\item \emph{Samplers.} Let $\Gamma(x,y)=\pseudorandomSamp(x)_y$. Then $\pseudorandomSamp$ is a $(\delta,\eps)$ averaging sampler if and only if $|\pseudorandomList_\Gamma(f,\mu(f)+\eps)|\le\delta N$ for every $f\colon[M]\to[0,1]$, and a Boolean one if and only if this holds for indicator functions.
\item \emph{Extractors.} Let $N=2^n$, $D=2^d$, $M=2^m$, and $K=2^k$. If $\Gamma$ is a $(k,\eps)$ extractor, then $|\pseudorandomList_\Gamma(f,\mu(f)+\eps)|<K$ for every $f\colon[M]\to[0,1]$. Conversely, if $|\pseudorandomList_\Gamma(T,\mu(T)+\eps)|\le K$ for every $T\subseteq[M]$, then $\Gamma$ is a $(k+\log(1/\eps),2\eps)$ extractor.
\end{enumerate}
\end{proposition}

\begin{proof}
(i) If some $S$ of size $K$ has $|\Gamma(S)|<AK$, then $T=\Gamma(S)$ is a test with $|T|<AK$ and $S\subseteq\pseudorandomList_\Gamma(T,1)$. Conversely, if $|T|<AK$ and $\pseudorandomList_\Gamma(T,1)$ contains a set $S$ of size $K$, then $\Gamma(S)\subseteq T$ has fewer than $AK$ elements.

(ii) The pair $(y,\pseudorandomEnc(x)_y)$ lies in $T_r$ exactly when $\pseudorandomEnc(x)_y=r_y$; average over $y$. (iii) An input $x$ lies in $\pseudorandomList_\Gamma(f,\mu(f)+\eps)$ exactly when the sampler's estimate exceeds $\mu(f)+\eps$, so the statement is the definition, read as a count.

(iv) Suppose that $\Gamma$ is a $(k,\eps)$ extractor and that $\cL=\pseudorandomList_\Gamma(f,\mu(f)+\eps)$ has at least $K$ elements. The uniform distribution $X$ on $\cL$ is a $k$-source, and $\E f(\Gamma(X,y))$, the average of $\Gamma_x(f)$ over $\cL$, exceeds $\mu(f)+\eps$, which contradicts \cref{lem:pseudorandom:tests}(i). For the converse, let $X$ have min-entropy at least $k+\log(1/\eps)$, let $T\subseteq[M]$, and let $\cL=\pseudorandomList_\Gamma(T,\mu(T)+\eps)$. Every $x\notin\cL$ has $\Gamma_x(T)\le\mu(T)+\eps$ and every point has probability at most $\eps/K$, so $\Pr[\Gamma(X,y)\in T]\le\Pr[X\in\cL]+\mu(T)+\eps\le\mu(T)+2\eps$, and \cref{lem:pseudorandom:tests}(i) turns this into $\Delta\le2\eps$.
\end{proof}

\begin{figure}[t]
\centering
\begin{tikzpicture}[x=1cm,y=0.62cm,
    lv/.style={circle,fill=seadeep,inner sep=0pt,minimum size=4.4pt},
    rv/.style={circle,fill=muted,inner sep=0pt,minimum size=4.4pt}]
  \fill[seamist,rounded corners=5pt] (5.62,1.95) rectangle (6.38,3.67);
  \node[font=\small,text=seadeep,anchor=west] at (6.5,2.81) {$T$};
  \fill[sand,rounded corners=5pt] (-0.38,2.55) rectangle (0.38,4.05);
  \draw[sandline,rounded corners=5pt] (-0.38,2.55) rectangle (0.38,4.05);
  \node[font=\small,anchor=east] at (-0.5,3.3) {$\pseudorandomList_\Gamma(T,\theta)$};
  \foreach \i in {0,...,7} \node[lv] (L\i) at (0,{0.55*\i}) {};
  \foreach \j in {0,...,9} \node[rv] (R\j) at (6,{0.43*\j}) {};
  \foreach \a/\b in {0/0,0/4,0/9,0/2,1/0,1/2,1/4,1/9,2/1,2/3,2/8,2/9,3/0,3/4,3/6,3/9,4/2,4/3,4/5,4/7,5/5,5/6,5/7,5/1,7/6,7/7,7/8,7/3}
    \draw[muted!45,line width=0.35pt] (L\a)--(R\b);
  \foreach \b in {5,6,8,1} \draw[sea,line width=0.9pt] (L6)--(R\b);
  \node[lv,fill=sea,minimum size=5.4pt] at (L6) {};
  \node[font=\small,text=sea,anchor=south] at (2.6,3.45) {$\Gamma(x,y)$, $y\in[D]$};
  \node[font=\small,anchor=east,text=sea] at (-0.12,4.75) {$x$};
  \draw[sea,line width=0.4pt,-{Stealth[length=4pt]}] (-0.14,4.62) -- (-0.05,3.42);
  \node[font=\small] at (0,-0.55) {$[N]$};
  \node[font=\small] at (6,-0.55) {$[M]$};
\end{tikzpicture}
\caption{One bipartite graph $\Gamma\colon[N]\times[D]\to[M]$, here with $D=4$. A left vertex $x$ is the randomness we pay for, and the label $y$ is enumerated; $\Gamma(x,y)$ is a neighbor, a sample, a coordinate with its symbol, or an output string. A test $T$ fools the left vertices whose edges land in it too often, and these form the list. The highlighted vertex sends three of its four edges into $T$, so it lies in $\pseudorandomList_\Gamma(T,\theta)$ for every $\theta<3/4$.}
\label{fig:pseudorandom:graph}
\end{figure}

\begin{table}[t]
\begin{fullwidth}
\centering\small
\renewcommand{\arraystretch}{1.2}
\begin{tabular}{@{}lllll@{}}
\toprule
Object & $\Gamma(x,y)$ & Tests & Threshold & Requirement on the list\\
\midrule
$({=}K,A)$ vertex expander & $y$th neighbor of $x$ & $|T|<AK$ & $1$ & fewer than $K$\\
$(\rho,L)$ list-decodable code & $(y,\pseudorandomEnc(x)_y)$ & $T_r$ for $r\in[q]^D$ & $1-\rho$ & at most $L$\\
$(\delta,\eps)$ averaging sampler & $\pseudorandomSamp(x)_y$ & all $f\colon[M]\to[0,1]$ & $\mu(f)+\eps$ & at most $\delta N$\\
$(k,\eps)$ extractor & $\pseudorandomExt(x,y)$ & all $T\subseteq[M]$ & $\mu(T)+\eps$ & fewer than $2^k$ (necessary)\\
lossless condenser & $\operatorname{Con}(x,y)$ & $|T|<(1-\eps)2^kD$ & $1$ & fewer than $2^k$\\
\bottomrule
\end{tabular}
\caption{The dictionary, after \textcite[Table~8.1]{Vadhan2012}. For extractors the bound ``at most $2^k$'' suffices up to $\log(1/\eps)$ bits of entropy and a factor $2$ in the error (\cref{prop:pseudorandom:dictionary}(iv)); lossless condensers appear in \cref{ex:pseudorandom:condenser}.}
\label{tab:pseudorandom:dictionary}
\end{fullwidth}
\end{table}

\Cref{fig:pseudorandom:graph} shows the picture and \cref{tab:pseudorandom:dictionary} collects the entries. Three dials distinguish the objects: the tests, the threshold, and the allowed size of the list. Their typical parameters are far apart: expanders want $D$ constant and $M=N$, codes want $q$ small and $D$ polynomial in $n=\log N$, and extractors and samplers want $\log D=O(\log(n/\eps))$ and $M$ much smaller than $N$ \parencite[\S8.1]{Vadhan2012}.

What did one have to notice? Not a deep theorem but a pattern of quantifiers: in each definition, \emph{for every test, few inputs are fooled}. Samplers call the count an error probability, extractors a min-entropy threshold, codes a list size. The dictionary grew one entry at a time:\sidenote{These attributions follow the chapter notes of \textcite[\S5.6 and \S6.5]{Vadhan2012}.} \textcite{Sipser1988} matched hitting samplers with expanders, \textcite{NisanZuckerman1996} saw extractors as expanding graphs, \textcite{Zuckerman1997} identified extractors with samplers, \textcite{Trevisan2001} turned generators and codes into extractors, \textcite{TaShmaZuckerman2004} crystallized the list-decoding view, and \textcite{GuruswamiUmansVadhan2009} brought in vertex expansion.

\section{After: one inequality and one union bound}\label{sec:pseudorandom:after1}

\subsection{One Cauchy--Schwarz inequality}\label{sec:pseudorandom:collision}

The last two steps of the classical proof of the leftover hash lemma, read with one test in mind, give a lemma about a single random variable $Z$; in applications $Z$ is a random neighbor of a random element of a set.

\begin{lemma}[Collision lemma]\label{lem:pseudorandom:collision}
Let $Z$ be a random variable on $[M]$ and let $T\subseteq[M]$. Then
\[
  \bigl|\Pr[Z\in T]-\mu(T)\bigr|\le\sqrt{\bigl(M\pseudorandomCP(Z)-1\bigr)\,\mu(T)\bigl(1-\mu(T)\bigr)}.
\]
\end{lemma}

\begin{proof}
Write $p_z=\Pr[Z=z]$. Because the numbers $p_z-1/M$ sum to zero, we may subtract the constant $\mu(T)$ from the indicator of $T$:
\[
  \Pr[Z\in T]-\mu(T)=\sum_z\Bigl(p_z-\frac1M\Bigr)\bigl(\ind_T(z)-\mu(T)\bigr).
\]
By the Cauchy--Schwarz inequality the absolute value of the sum is at most the product of the square roots of $\sum_z(p_z-1/M)^2=\pseudorandomCP(Z)-1/M$ and $\sum_z(\ind_T(z)-\mu(T))^2=M\mu(T)(1-\mu(T))$.
\end{proof}

The lemma reduces ``is a random neighbor of a random element of $S$ nearly uniform?'' to a count of collisions $\Gamma(x,y)=\Gamma(x',y')$ with $x,x'\in S$, and each community had its own way to count them.

\begin{proof}[Proof of the mixing lemma (\cref{thm:pseudorandom:mixing})]
Both sides vanish for $S=\emptyset$, so let $s=|S|\ge1$ and let $Z=\Gamma(x,y)$ with $x$ uniform on $S$ and $y$ uniform on $[D]$. Let $\pi$ and $u$ be the uniform distributions on $S$ and on $[N]$, written as vectors. Because $P$ is symmetric, $Z$ has distribution $P\pi$. Because $Pu=u$ and $\pi-u\perp\mathbf1$, the definition of $\lambda$ gives $\norm{P\pi-u}^2\le\lambda^2\norm{\pi-u}^2=\lambda^2(1/s-1/N)$. Since $P\pi-u\perp u$, we have $\pseudorandomCP(Z)=\norm{P\pi}^2=1/N+\norm{P\pi-u}^2$, hence $N\pseudorandomCP(Z)-1\le\lambda^2(1-\sigma)/\sigma$. \Cref{lem:pseudorandom:collision} gives $|\Pr[Z\in T]-\tau|\le\lambda\sqrt{(1-\sigma)\tau(1-\tau)/\sigma}$; multiply by $\sigma$ and use $\sigma\Pr[Z\in T]=e(S,T)/(ND)$.
\end{proof}

\begin{proof}[Proof of the Johnson bound (\cref{thm:pseudorandom:johnson})]
Use the graph $\Gamma(x,y)=(y,\pseudorandomEnc(x)_y)$ of \cref{prop:pseudorandom:dictionary}(ii), with $M=Dq$. Let $S$ be the set of messages whose codewords agree with $r$ on more than a $1/q+\eps$ fraction of coordinates, suppose $s=|S|\ge1$, and let $Z=\Gamma(x,y)$ with $x$ uniform on $S$ and $y$ uniform on $[D]$. Two independent copies of $Z$ coincide when they use the same coordinate, which has probability $1/D$, and carry the same symbol there. Of the $s^2$ ordered pairs of messages, $s$ are equal and the others agree on at most a $1/q+\eta$ fraction of coordinates, so
\[
  M\pseudorandomCP(Z)\le\frac q{s^2}\Bigl(s+s(s-1)\bigl(\tfrac1q+\eta\bigr)\Bigr)\le1+\frac{q-1}s+q\eta .
\]
Since $\Pr[Z\in T_r]$ is the average agreement over $S$, it exceeds $\mu(T_r)+\eps$, and \cref{lem:pseudorandom:collision} with $\mu(T_r)(1-\mu(T_r))=(q-1)/q^2$ yields $\eps^2<\bigl((q-1)/s+q\eta\bigr)(q-1)/q^2$. This rearranges to the stated bound on $s$.
\end{proof}

\begin{proof}[Second proof of the leftover hash lemma (\cref{thm:pseudorandom:lhl})]
The collision bound $\pseudorandomCP(Z)\le(2^{-k}+1/M)/|\cH|$ for $Z=(h,h(X))$ gives $|\cH|M\pseudorandomCP(Z)-1\le M2^{-k}\le\eps^2$. Since $\mu(1-\mu)\le1/4$, \cref{lem:pseudorandom:collision} bounds $|\Pr[Z\in T]-\mu(T)|$ by $\eps/2$ for every test $T\subseteq\cH\times[M]$.
\end{proof}

Graph theorists count collisions with an eigenvalue, coding theorists with the distance, and cryptographers with universality; the rest is one inequality, as \textcite{TaShmaZuckerman2004} had noticed for the last two.

\subsection{One union bound}\label{sec:pseudorandom:union}

The random constructions also become one, and the price of a test is a relative entropy. For $\theta\in[0,1]$ and $\mu\in(0,1)$ let $\pseudorandomKL(\theta\Vert\mu)\defeq\theta\ln(\theta/\mu)+(1-\theta)\ln((1-\theta)/(1-\mu))$, with $0\ln0=0$.

\begin{lemma}[Chernoff's bound]\label{lem:pseudorandom:chernoff}
Let $B_1,\dots,B_t$ be independent random variables with values in $\{0,1\}$ whose means average to $\bar\mu\in(0,1)$, and let $\theta\in[\bar\mu,1]$. Then $\Pr[\sum_iB_i\ge\theta t]\le\exp(-t\pseudorandomKL(\theta\Vert\bar\mu))$. Moreover $\pseudorandomKL(\theta\Vert\bar\mu)\ge2(\theta-\bar\mu)^2$ and $\pseudorandomKL(1\Vert\bar\mu)=\ln(1/\bar\mu)$.
\end{lemma}

\begin{proof}
For $\theta=1$ the probability is the product of the means, at most $\bar\mu^t$ by the inequality of arithmetic and geometric means. For $\theta<1$ and $\beta\ge0$, Markov's inequality for $e^{\beta\sum_iB_i}$ and the same inequality give the bound $(e^{-\beta\theta}(1-\bar\mu+\bar\mu e^\beta))^t$, which equals the claim for $e^\beta=\theta(1-\bar\mu)/(\bar\mu(1-\theta))$. Finally, $\pseudorandomKL(\theta\Vert\bar\mu)-2(\theta-\bar\mu)^2$ vanishes with its derivative at $\theta=\bar\mu$ and has second derivative $1/(\theta(1-\theta))-4\ge0$.
\end{proof}

\begin{theorem}[One union bound]\label{thm:pseudorandom:union}
Let $K\le N$, let $\cT$ be a family of subsets $T$ of $[M]$ with $0<|T|<M$, and for each $T\in\cT$ let $\theta(T)\in[\mu(T),1]$. Let $\Gamma\colon[N]\times[D]\to[M]$ be random in one of two ways: either its $ND$ values are independent and uniform on $[M]$, or $M=Dq$ and $\Gamma(x,y)=(y,F(x,y))$ with the values $F(x,y)$ independent and uniform on $[q]$. The probability that some $S\subseteq[N]$ with $|S|=K$ and some $T\in\cT$ satisfy $\Pr_{x\in S,y}[\Gamma(x,y)\in T]\ge\theta(T)$ is at most
\[
  \binom NK\sum_{T\in\cT}\exp\bigl(-KD\pseudorandomKL(\theta(T)\Vert\mu(T))\bigr).
\]
Outside this event, $|\pseudorandomList_\Gamma(T,\theta(T))|<K$ for every $T\in\cT$.
\end{theorem}

\begin{proof}
Fix $S$ and $T$. The $KD$ indicators of the events $\Gamma(x,y)\in T$, for $x\in S$ and $y\in[D]$, are independent. In the first model each has mean $\mu(T)$; in the second, the mean is $|T_y|/q$ with $T_y=\{a:(y,a)\in T\}$, and these means average to $\mu(T)$. \Cref{lem:pseudorandom:chernoff} bounds the probability for the pair $(S,T)$, and a union bound over the pairs proves the first claim. For the second, a list with at least $K$ elements contains a set $S$ of size $K$, over which $\Gamma_x(T)$ averages at least $\theta(T)$.
\end{proof}

\Textcite[Theorem~8.2]{Vadhan2012} gives a sharper nonconstructive unified object of this kind. Threshold $1$ and tests of size $(D-2)K$ give Pinsker's theorem (\cref{ex:pseudorandom:randomexp}); extractors and codes come next.

\begin{corollary}[Random extractors]\label{cor:pseudorandom:randomext}
Let $0\le k\le n$ be integers and $\eps\in(0,1)$, and put $d=\ceil{\log(n-k+1)+2\log(1/\eps)}$ and $m=\floor{k+d-2\log(1/\eps)}$. Some $\Gamma\colon\{0,1\}^n\times\{0,1\}^d\to\{0,1\}^m$ is a $(k,\eps)$ extractor.
\end{corollary}

\begin{proof}
Use the first model with $K=2^k$, all tests with $0<\mu(T)\le1-\eps$, and $\theta(T)=\mu(T)+\eps$; other tests cannot exceed their density by $\eps$. Each term of the sum is at most $e^{-2KD\eps^2}$ by \cref{lem:pseudorandom:chernoff}, there are fewer than $2^M$ tests, and $\binom NK\le(eN/K)^K$, so the bad event has probability less than $\exp(K\ln(eN/K)+M\ln2-2KD\eps^2)$. The choices of $d$ and $m$ give $D\eps^2\ge n-k+1\ge\ln(eN/K)$ and $M\le KD\eps^2$, so the exponent is negative. Hence, with positive probability, $\Pr_{x\in S,y}[\Gamma(x,y)\in T]<\mu(T)+\eps$ for every $S$ of size $K$ and every $T$, and \cref{lem:pseudorandom:tests}(ii) finishes the proof.
\end{proof}

In the relevant ranges, both the seed length and the entropy loss $k+d-m\approx2\log(1/\eps)$ are optimal up to additive constants \parencite{RadhakrishnanTaShma2000}.

\begin{proof}[Proof of \cref{thm:pseudorandom:elias} (random codes)]
If $N\le L$ there is nothing to prove. Otherwise use the second model on the code's graph, with $K=L+1$, the $q^D$ tests $T_r$ of density $1/q$, and the threshold $1-\rho>1/q$; a word with $L+1$ messages of agreement above $1-\rho$ would give a set $S$ whose average agreement exceeds the threshold. Expanding the definitions,
\[
  \pseudorandomKL\Bigl(1-\rho\Bigm\Vert\frac1q\Bigr)=(1-\rho)\ln\bigl((1-\rho)q\bigr)+\rho\ln\frac{\rho q}{q-1}=\bigl(1-H_q(\rho)\bigr)\ln q,
\]
so with $N\le q^{RD}$ the bad event has probability at most
\[
  N^{L+1}q^Dq^{-(L+1)D(1-H_q(\rho))}\le q^{(L+1)D\,(R+1/(L+1)-1+H_q(\rho))}<1.\qedhere
\]
\end{proof}

\subsection{Translations}\label{sec:pseudorandom:translations}

The dictionary turns a statement about one object into a statement about another, at a stated exchange rate.

\begin{theorem}[Extractors are samplers; \textcite{Zuckerman1997}]\label{thm:pseudorandom:zuckerman}
Let $\Gamma\colon[N]\times[D]\to[M]$ with $N=2^n$, $D=2^d$, and $M=2^m$, and let $\pseudorandomSamp(x)\defeq(\Gamma(x,y))_{y\in[D]}$.
\begin{enumerate}
\item If $\Gamma$ is a $(k,\eps)$ extractor, then $\pseudorandomSamp$ is a $(2^{k-n},\eps)$ averaging sampler.
\item If $\pseudorandomSamp$ is a $(\delta,\eps)$ Boolean averaging sampler, then $\Gamma$ is an $(n-\log(1/\delta)+\log(1/\eps),2\eps)$ extractor.
\end{enumerate}
\end{theorem}

\begin{proof}
Each part combines \cref{prop:pseudorandom:dictionary}(iii) and (iv), with $K=2^k$ in the first part and $K=\delta N$ in the second.
\end{proof}

The sampler's random bits are the source length $n$, its queries are the $2^d$ seeds, and its confidence $\log(1/\delta)$ is the entropy deficiency $n-k$; the translation is exact one way and costs $\log(1/\eps)$ bits and a factor $2$ the other.

For instance, \cref{cor:pseudorandom:randomext} with $k=m$ and $n=m+t$, keeping the first $m$ output bits, gives by part~(i) a $(2^{-t},\eps)$ averaging sampler for functions on $\{0,1\}^m$ with $m+t$ random bits and fewer than $2(t+1)/\eps^2$ queries: the economy of pairwise independence in bits and of independent samples in queries, at once. Explicit extractors give explicit samplers with $O(m+\log(1/\delta))$ random bits and $\poly(1/\eps,\log(1/\delta))$ queries \parencite[Table~8.1]{Vadhan2012}.\index{deterministic amplification}

\begin{corollary}[Amplification and weak sources]\label{cor:pseudorandom:amplify}
Let $A(w;r)$ use $r\in\{0,1\}^m$ and err with probability at most $1/3$ on each input $w$. Let $\pseudorandomExt\colon\{0,1\}^n\times\{0,1\}^d\to\{0,1\}^m$ be a $(k,1/10)$ extractor, and let $A'(w;x)$ be the majority vote of $A(w;\pseudorandomExt(x,y))$ over all $y\in\{0,1\}^d$. For every $t\ge0$ and every random variable $X$ on $\{0,1\}^n$ with $\pseudorandomHmin(X)\ge k+t$, the algorithm $A'(w;X)$ errs with probability less than $2^{-t}$.
\end{corollary}

\begin{proof}
Fix $w$ and let $B$ be the set of random strings on which $A$ errs, so $\mu(B)\le1/3$. If the vote on $x$ is wrong, counting ties as wrong, then at least half of the strings $\pseudorandomExt(x,y)$ lie in $B$, so $\pseudorandomExt_x(B)\ge1/2>\mu(B)+1/10$ and $x\in\pseudorandomList_{\pseudorandomExt}(B,\mu(B)+1/10)$. This list has fewer than $2^k$ elements by \cref{prop:pseudorandom:dictionary}(iv), and each has probability at most $2^{-k-t}$ under $X$.
\end{proof}

For uniform $x$ this is error reduction: with the extractor of \cref{cor:pseudorandom:randomext} ($k=m$, $\eps=1/10$), error $2^{-t}$ costs $m+t$ random bits and fewer than $200(t+1)$ runs. For a weak source it is the simulation of randomized algorithms with weak sources, earlier proved separately for class after class \parencite[\S6.5]{Vadhan2012}. All of the difficulty has moved into the construction of the extractor.

Codes translate likewise: by \cref{prop:pseudorandom:dictionary}(ii) and (iv) a strong $(k,\eps)$ extractor with output alphabet $[q]$ is a $(1-1/q-\eps,2^k)$ list-decodable code, and conversely, since the tests $T_r$ are the worst ones for one output bit, a binary $(1/2-\eps,L)$ list-decodable code is a strong $(\log L+\log(1/\eps),2\eps)$ extractor, though for alphabet $[q]$ the error becomes $q\eps$ \parencite[Proposition~6.25]{Vadhan2012}.

\section{After: constructions cross the bridges}\label{sec:pseudorandom:after2}

A common language would be a mere renaming if every construction still had to be invented separately for every object. It is not, because constructions moved along the dictionary in both directions.

\subsection{From pseudorandom generators to extractors}\label{sec:pseudorandom:trevisan}

\textcite{NisanWigderson1994} built pseudorandom generators by evaluating a hard function on many nearly disjoint pieces of a short seed. \Textcite{Trevisan2001} observed that the construction is an extractor when the hard function is replaced by the encoding of a weak random sample under a list-decodable binary code.\index{Trevisan's extractor} The pieces come from a design.

\begin{definition}[Design]\label{def:pseudorandom:design}
Sets $S_1,\dots,S_m\subseteq[d]$ form an \emph{$(\ell,a)$ design} if $|S_i|=\ell$ for every $i$ and $|S_i\cap S_j|\le a$ for all $i\ne j$. For $y\in\{0,1\}^d$ and $S\subseteq[d]$, the string $y|_S$ lists the coordinates of $y$ indexed by $S$, in increasing order.
\end{definition}

Designs\index{combinatorial design} come from polynomials. If $p$ is a power of $2$ and $B\subseteq\F_p$ has $\ell\le p$ elements, then the sets $\{(\beta,g(\beta)):\beta\in B\}$, one for each polynomial $g$ over $\F_p$ of degree at most $a$, form an $(\ell,a)$ design of $p^{a+1}$ sets in a universe of size $d=p^2$, because two distinct such polynomials agree in at most $a$ points. For a code $C\colon\{0,1\}^n\to\{0,1\}^{2^\ell}$ write $C_x\colon\{0,1\}^\ell\to\{0,1\}$ for the codeword of $x$, read as a function, and set
\[
  \pseudorandomExt(x,y)\defeq\bigl(C_x(y|_{S_1}),\dots,C_x(y|_{S_m})\bigr)\in\{0,1\}^m .
\]

\begin{theorem}[\textcite{Trevisan2001}]\label{thm:pseudorandom:trevisan}
Let $S_1,\dots,S_m\subseteq[d]$ be an $(\ell,a)$ design, let $\eps>0$, and let $C\colon\{0,1\}^n\to\{0,1\}^{2^\ell}$ be a code such that, for every $r\colon\{0,1\}^\ell\to\{0,1\}$, at most $L$ messages $x$ satisfy $\pseudorandomAgr(C_x,r)>1/2+\eps/m$. Then $|\pseudorandomList_{\pseudorandomExt}(T,\mu(T)+\eps)|\le Lm2^{m+(m-1)2^a}$ for every $T\subseteq\{0,1\}^m$, and $\pseudorandomExt$ is a $(k,2\eps)$ extractor for $k=\log L+\log m+m+(m-1)2^a+\log(1/\eps)$.
\end{theorem}

\begin{proof}
The plan is to show that each $x$ in the list has a short description relative to $T$: an index $i\le m$ and at most $m+(m-1)2^a$ further bits determine a word $r$ that agrees with $C_x$ on more than a $1/2+\eps/m$ fraction of inputs, and at most $L$ messages are consistent with $r$.

\emph{Hybrids.} Fix $T$ and $x\in\pseudorandomList(T,\mu(T)+\eps)$, write $b_j=C_x(y|_{S_j})$ for uniform $y\in\{0,1\}^d$, let $u\in\{0,1\}^m$ be uniform and independent of $y$, and for $0\le i\le m$ and $c\in\{0,1\}$ put
\[
  H_i=(b_1,\dots,b_i,u_{i+1},\dots,u_m),\qquad H_i^c=(b_1,\dots,b_{i-1},c,u_{i+1},\dots,u_m).
\]
Let $h_i=\Pr[H_i\in T]$. Then $h_0=\mu(T)$ and $h_m=\pseudorandomExt_x(T)>\mu(T)+\eps$, so some $i$ has $h_i-h_{i-1}>\eps/m$.

\emph{Prediction.} Let $c$ be a further uniform bit, and let $P$ output $c$ if $H_i^c\in T$ and $1-c$ otherwise. When $c=b_i$ the output is right exactly when $H_i\in T$, and when $c\ne b_i$ exactly when $H_i^{1-b_i}\notin T$. Since $h_{i-1}$ is the average of $\Pr[H_i^0\in T]$ and $\Pr[H_i^1\in T]$, averaging over $c$ gives $\Pr[P=b_i]=1/2+h_i-h_{i-1}>1/2+\eps/m$.

\emph{Fixing.} Split $y$ into $z=y|_{S_i}$ and the remaining coordinates $w$. By averaging, some fixed $w^*$, $u^*_{i+1},\dots,u^*_m$, and $c^*$ keep the success probability of $P$ above $1/2+\eps/m$ when only $z$ is random. With $w=w^*$, each $b_j$ with $j<i$ depends on $z$ only through the at most $a$ coordinates in $S_i\cap S_j$, so it is a function $g_j(z)$ given by a table of at most $2^a$ bits.

\emph{Counting.} Once $T$ is fixed, the word $r(z)\defeq P(z;w^*,u^*,c^*)$ is determined by $i$, by the $m-i+1$ bits $u^*$ and $c^*$, and by the tables of $g_1,\dots,g_{i-1}$. For each $i$ there are at most $2^{m-i+1+(i-1)2^a}\le2^{m+(m-1)2^a}$ such descriptions, hence at most $m2^{m+(m-1)2^a}$ words in all, and each $x$ in the list is one of at most $L$ messages for its word. The last claim is \cref{prop:pseudorandom:dictionary}(iv) with $K=Lm2^{m+(m-1)2^a}$.
\end{proof}

This is the encoding argument of \cref{ch:compression}: if the extractor fails on $x$, then $x$ has a description of about $m2^a+\log L$ bits relative to the test, and a source of higher min-entropy rarely produces such strings.\index{reconstruction paradigm} The code needs no decoder, only a bound on lists, and the hard function of Nisan and Wigderson has become a source of high min-entropy, because a black-box proof of pseudorandomness shows exactly that a function fooled by a test has a short description \parencite[\S7.7]{Vadhan2012}.

\begin{corollary}[An explicit extractor with a polylogarithmic seed]\label{cor:pseudorandom:explicit}
For every $\eps\in(0,1)$ and all integers $n\ge k\ge\max\{128,4+6\log(1/\eps)\}$ there is a $(k,\eps)$ extractor $\pseudorandomExt\colon\{0,1\}^n\times\{0,1\}^d\to\{0,1\}^{\floor{\sqrt k}}$ with $d=O(\log^2(n/\eps))$, computable in time polynomial in $n$ and $\log(1/\eps)$.
\end{corollary}

\begin{proofsketch}
Apply \cref{thm:pseudorandom:trevisan} with error parameter $\eps/2$ and $m=\floor{\sqrt k}$. For the code, read $x$ as a polynomial $f_x$ of degree less than $\ceil{n/t}$ over $\F_{2^t}$, with $2^t\ge2nm^2/\eps^2$, and let $C_x(\beta,v)$ be the inner product over $\F_2$ of $f_x(\beta)$ with $v\in\{0,1\}^t$, a Reed--Solomon code followed by a Hadamard code with $\ell=2t$. Distinct polynomials agree at fewer than $n/t$ points and distinct inner-product codewords agree on exactly half of the coordinates, so pairwise agreements are at most $1/2+\eps^2/(4m^2)$, and \cref{thm:pseudorandom:johnson} gives lists $L<2m^2/\eps^2$ at agreement $1/2+\eps/(2m)$. For the design take $a=\floor{\frac12\log k}-2$ and $p$ the least power of $2$ that is at least $\max\{\ell,16\}$; then $p^{a+1}\ge k^2/256\ge m$ and $d=p^2=O(\log^2(n/\eps))$. Since $2^a\le\sqrt k/4$, the min-entropy that the theorem requires is less than $2+3\log(1/\eps)+3\log m+m+k/4\le k$ in the stated range. All arithmetic is in fields of size polynomial in $n/\eps$, whose descriptions can be found in polynomial time \parencite[\S5.4]{Vadhan2012}. Only routine arithmetic is omitted; \textcite[Theorem~7.73]{Vadhan2012} gives a sharper version.
\end{proofsketch}

With designs that have $d=O(\ell^2/a)$ for $a=\gamma\log m$ \parencite[Lemma~7.22]{Vadhan2012}, the seed becomes $O(\log^2(n/\eps)/\log k)$ and the output $k^{1-\gamma}$ for every constant $\gamma>0$ \parencite[Theorem~7.73]{Vadhan2012}. Compare the length of this proof with the aside of \cref{sec:pseudorandom:extractors}.

\subsection{From codes to expanders, and from extractors to graphs}\label{sec:pseudorandom:guv}

Parvaresh and Vardy found, in 2005, codes that can be list-decoded in polynomial time beyond the radius $1-\sqrt R$ of Guruswami and Sudan,\sidenote{See \textcite[\S5.6]{Vadhan2012}.} and \textcite{GuruswamiUmansVadhan2009} read them through \cref{prop:pseudorandom:dictionary}(ii) as graphs.

\begin{theorem}[\textcite{GuruswamiUmansVadhan2009}]\label{thm:pseudorandom:guv}
Let $q$ be a prime power, let $n,m,h\ge1$ be integers, and let $E\in\F_q[Y]$ be irreducible of degree $n$. For $f\in\F_q[Y]$ of degree less than $n$ let $f_i\defeq f^{h^i}\bmod E$ for $0\le i<m$. Then $\Gamma(f,y)\defeq(y,f_0(y),\dots,f_{m-1}(y))$ defines an $(h^m,\,q-nhm)$ vertex expander $\Gamma\colon\F_q^n\times\F_q\to\F_q^{m+1}$.
\end{theorem}

The proof \parencite[Theorem~5.34]{Vadhan2012} is the Parvaresh--Vardy decoding argument with the graph $T_r$ of a received word replaced by a small test $T$, as \cref{prop:pseudorandom:dictionary}(i) suggests: some nonzero low-degree polynomial $Q(Y,Z_0,\dots,Z_{m-1})$ vanishes on $T$, because $|T|<AK$ is less than the number of monomials; then $Q(Y,f_0,\dots,f_{m-1})=0$ for every $f\in\pseudorandomList(T,1)$, and such $f$ are among the few roots of $Q(Y,Z,Z^h,\dots,Z^{h^{m-1}})$ over $\F_q[Y]/(E)$. It is the polynomial method of \cref{ch:dimension}, applied to a graph. Choosing parameters gives explicit $(K,(1-\eps)D)$ expanders with degree $D=O((\log N)(\log K)/\eps)^{1+1/\alpha}$ and at most $D^2K^{1+\alpha}$ right vertices, for every constant $\alpha>0$ \parencite[Theorem~5.35]{Vadhan2012}. They are lossless condensers\index{lossless condenser} (\cref{ex:pseudorandom:condenser}), and composed with simpler extractors they give extractors optimal up to constant factors.

\begin{theorem}[\textcite{GuruswamiUmansVadhan2009}]\label{thm:pseudorandom:guvext}
For every constant $\gamma>0$, all integers $n\ge k$, and every $\eps>0$, there is an explicit $(k,\eps)$ extractor $\pseudorandomExt\colon\{0,1\}^n\times\{0,1\}^d\to\{0,1\}^m$ with $d=O(\log(n/\eps))$ and $m\ge(1-\gamma)k$.
\end{theorem}

Constructions also moved from extractors to graphs. \textcite{ReingoldVadhanWigderson2002} analyzed their zig-zag product\index{zig-zag product} by following how entropy moves between the two halves of a vertex, as one analyzes extractors; the block-source composition of \cref{sec:pseudorandom:extractors} is its twin \parencite[\S6.3.5]{Vadhan2012}. For a $D_1$-regular $G$ on $N_1$ vertices and a $D_2$-regular $H$ on $D_1$ vertices, $G\pseudorandomZigzag H$ replaces each vertex of $G$ by a cloud of $D_1$ vertices; a step is a step in $H$, a step in $G$ between clouds, and another step in $H$.

\begin{theorem}[Zig-zag product; \textcite{ReingoldVadhanWigderson2002}]\label{thm:pseudorandom:zigzag}
If $\lambda(G)\le\lambda_1$ and $\lambda(H)\le\lambda_2$, then $G\pseudorandomZigzag H$ is a $D_2^2$-regular graph on $N_1D_1$ vertices with $\lambda(G\pseudorandomZigzag H)\le\lambda_1+2\lambda_2$.
\end{theorem}

This form of the bound follows from \textcite[Theorem~4.35]{Vadhan2012}. Alternately squaring and zig-zagging with a fixed small expander gives explicit constant-degree expanders by an elementary recursion, and the new reach was in complexity theory: \textcite{Reingold2008} proved that connectivity in undirected graphs can be decided deterministically in logarithmic space. His algorithm, announced in 2005, uses the zig-zag product to turn each component of the input graph into an expander of logarithmic diameter within logarithmic space, and then tries all short walks \parencite[\S4.4]{Vadhan2012}.\sidenote{In the language of complexity classes, $\mathrm{SL}=\mathrm{L}$.}

\section{Costs and limits}\label{sec:pseudorandom:costs}

\paragraph{Where did the difficulty go?} To three places. The first is the explicit constructions, which the dictionary transports but does not create: Trevisan's extractor needs a list-decodable code and a design, the expanders of Guruswami, Umans, and Vadhan need the algebra of Parvaresh--Vardy codes, and Ramanujan graphs need number theory. The second is the exchange rates: samplers give extractors only after losing $\log(1/\eps)$ bits and doubling the error, codes give extractors with the error multiplied by the alphabet size, and vertex expansion tolerates no constant-factor loss in a list, since that is the same loss in expansion \parencite[\S5.3]{Vadhan2012}. The third is one inequality and one union bound, each paid once. The bargain is that a construction of one object is at once a construction of the others wherever their parameters overlap, as when a generator became an extractor and codes became lossless expanders.

\paragraph{What did not dissolve.} The dictionary unifies definitions, not regimes, and a translation helps only where the parameter ranges meet. Lossless expansion with constant degree needed a new object, the randomness conductors of \textcite{CapalboReingoldVadhanWigderson2002}. \Textcite{Vadhan2012} lists as open, among others, explicit extractors with seed $O(\log n)$ and output $k+d-O(1)$ for constant $\eps$ (Open Problem~6.41), explicit binary codes that achieve list-decoding capacity (Open Problem~5.26), and one explicit $\Gamma$ that matches the union bound of \cref{thm:pseudorandom:union} for all tests at once (Open Problem~8.5). Extraction without a seed lies outside the framework: the explicit two-source extractor of \textcite{ChattopadhyayZuckerman2019}, for min-entropy $\log^Cn$ where the previous best needed about $0.499n$, required new ideas. The framework is also information-theoretic: a generator that fools all small circuits would yield a hard function, so building one means proving circuit lower bounds, about which the dictionary is silent \parencite[\S7]{Vadhan2012}. Finally, unification costs constants: the union over tests proves Pinsker's theorem with a density exponentially small in the degree, where counting repeats gives a polynomial one (\cref{ex:pseudorandom:randomexp}).

\paragraph{Other seas.} Vadhan's theory also contains black-box generators and hardness amplifiers, which are locally list-decodable codes. In the algebraic world the move recurs on its own: small-bias distributions on $\{0,1\}^n$, linear codes whose nonzero codewords have weight near $1/2$, and expanding Cayley graphs on $\F_2^n$ are one object, seen through the Fourier transform of \cref{ch:fourier} \parencite[\S8.2.5]{Vadhan2012}. Restricting the tests changes the objects, as for the $\eps$-approximations of \cref{ch:vc}.

\section{The move, distilled}\label{sec:pseudorandom:distilled}

\begin{move}{Unification}
\movepart{In one sentence} When several objects are each defined by ``for every test, few inputs are fooled'', put them on one bipartite graph, measure each by the size of one set, the inputs that a test fools, and move proofs and constructions along the resulting dictionary.
\movepart{Triggers} The same union bound or Cauchy--Schwarz step in papers from different fields; parameters that recur across literatures, such as $\log(1/\delta)/\eps^2$ queries and $2\log(1/\eps)$ bits of loss; constructions in one field quietly built from the objects of another.
\movepart{Recipe} (1) Find a common syntax: a function of one argument that is paid for and one that is enumerated. (2) Rewrite each property as a bound on a set of fooled inputs, recording the tests and the threshold. (3) Compute the exchange rate of every translation. (4) Prove the generic theorems once. (5) Transport constructions, after checking that the parameter regimes overlap.
\movepart{Anti-pattern} A dictionary that ignores regimes, such as calling an extractor of polylogarithmic degree an expander when the application needs constant degree. Equally fake is a common definition so loose that each instance must prove its properties from scratch.
\end{move}

\dissolutiontest{Mixing lemma, Johnson bound, and leftover hash lemma from one collision lemma; weak-source amplification in three lines; Trevisan's extractor in a page.}{Every extractor is a sampler, with no independence to analyze; Trevisan's construction needs list bounds, not decoders, and min-entropy, not hardness.}{Expanders, codes, samplers, extractors, condensers, and black-box generators are bounds on $|\pseudorandomList_\Gamma(T,\theta)|$.}{Optimal samplers; extractors from generators; lossless expanders from codes; zig-zag and $\mathrm{SL}=\mathrm{L}$.}

Among the unifications in this book this one is the most syntactic. Linear programming duality (\cref{ch:lpduality}) derived the minimax theorems from one theorem; here the objects were already one, and the work was to notice it, prove the generic facts once, and keep honest books of the exchange rates. Trevisan's reconstruction is the encoding move of \cref{ch:compression}, the collision lemma is a second-moment computation of the kind Parseval's identity performs in \cref{ch:fourier}, and the algebraic constructions are the polynomial method of \cref{ch:dimension}.

\begin{exercises}
\item \consolidation\label{ex:pseudorandom:disperser} A function $\Gamma\colon[N]\times[D]\to[M]$ is a \emph{$(k,\eps)$ disperser} if $|\Gamma(S)|\ge(1-\eps)M$ whenever $|S|=2^k$. (a) Show that this holds if and only if $|\pseudorandomList_\Gamma(T,1)|<2^k$ for every $T$ with $|T|<(1-\eps)M$. (b) Show that every $(k,\eps)$ extractor is a $(k,\eps)$ disperser, and that a $(k,1/3)$ disperser reduces one-sided error $1/2$ to less than $2^{k-\log N}$ (run the algorithm on every $\Gamma(x,y)$ and accept if some run accepts).

\item \consolidation\label{ex:pseudorandom:hadamard} The Hadamard code sends $x\in\{0,1\}^n$ to the word whose coordinate $v\in\{0,1\}^n$ is the inner product $\ip{x}{v}$ over $\F_2$. (a) Show that distinct codewords agree on exactly half of the coordinates. (b) Deduce from \cref{thm:pseudorandom:johnson} that fewer than $1/(4\eps^2)$ codewords agree with a given word on more than a $1/2+\eps$ fraction of coordinates. (c) Prove it with Parseval's identity (\cref{ch:fourier}) instead.

\item \consolidation\label{ex:pseudorandom:spectral} Let $\Gamma\colon[N]\times[D]\to[N]$, with $N=2^n$, be $D$-regular and undirected with $\lambda(\Gamma)\le\lambda$. (a) Use \cref{lem:pseudorandom:collision} to show that $\Gamma(X,y)$ is $\frac\lambda22^{(n-k)/2}$-close to uniform for every $k$-source $X$ and uniform $y$. (b) Deduce that if $\lambda\le2/\sqrt D$, then a seed of $n-k+2\log(1/\eps)$ bits suffices for a $(k,\eps)$ extractor. (Goldreich and Wigderson; see \textcite[\S6.2]{Vadhan2012}.)

\item \consolidation\label{ex:pseudorandom:randomexp} Use \cref{thm:pseudorandom:union} with threshold $1$ and the tests of size $(D-2)K$ to show that, for $D\ge3$ and $\alpha=1/(4e^{D-1}(D-2)^2)$, a random $\Gamma\colon[N]\times[D]\to[N]$ with independent uniform values is an $(\alpha N,D-2)$ vertex expander with probability greater than $2/3$. Compare $\alpha$ with the value $1/(e^3D^4)$ that the count of repeats in \textcite[Theorem~4.4]{Vadhan2012} gives. What does the union over tests waste?

\item \transfer\label{ex:pseudorandom:privacy} Privacy amplification. Alice and Bob share $X\in\{0,1\}^n$, an eavesdropper holds $E$ correlated with $X$, and they publish a uniform seed $Y$, independent of $(X,E)$, and keep $\pseudorandomExt(X,Y)$. (a) Show that if $\pseudorandomExt$ is a strong $(k,\eps)$ extractor and $\pseudorandomHmin(X\mid E=e)\ge k$ for every $e$, then $(Y,\pseudorandomExt(X,Y),E)$ is $\eps$-close to $(Y,U,E)$, with $U$ uniform and independent of $(Y,E)$. (b) Use \cref{prop:pseudorandom:dictionary}(iv) to show that $\sum_e\Pr[E=e]\max_x\Pr[X=x\mid E=e]\le\eps2^{-k}$ already gives distance at most $2\eps$.

\item \transfer\label{ex:pseudorandom:walks} Walks as samplers. Let $G$ be $D$-regular on $[N]$ with $\lambda(G)\le\lambda$, and let $\Gamma\colon([N]\times[D]^{t-1})\times[t]\to[N]$ send a start vertex, $t-1$ edge labels, and an index $i$ to the $i$th vertex of the walk they determine. (a) A walk stays inside a set of density $\mu$ with probability at most $(\mu+\lambda(1-\mu))^t$ \parencite[Theorem~4.17]{Vadhan2012}; rewrite this as a bound on a list. (b) Assume the Chernoff bound for walks, $\Pr[\frac1t\sum_if(V_i)\ge\mu(f)+\lambda+\eps]\le2e^{-c\eps^2t}$ with $c>0$ absolute \parencite[Theorem~4.22]{Vadhan2012}, whose first form is due to \textcite{Gillman1998}. Translate it into a sampler and, by \cref{thm:pseudorandom:zuckerman}, into an extractor; give its seed length and required min-entropy.

\item \transfer\hard\label{ex:pseudorandom:xor} Hardness amplification is decoding. Let $f\colon\{0,1\}^\ell\to\{0,1\}$, and let $H\subseteq\{0,1\}^\ell$ have density at least $\delta$ and be such that every circuit of size at most $s$ agrees with $f$ on at most a $1/2+\eps$ fraction of $H$; the hardcore lemma of \cref{ch:lpduality} provides such a set when $f$ is mildly hard for somewhat larger circuits. Let $f^{\oplus t}(x_1,\dots,x_t)=f(x_1)\oplus\dots\oplus f(x_t)$. (a) Show that every circuit of size at most $s-1$ agrees with $f^{\oplus t}$ on at most a $1/2+\eps+(1-\delta)^t/2$ fraction of inputs. (b) Explain why this is a reconstruction argument, as in \cref{thm:pseudorandom:trevisan}. \hint{condition on the first index $i$ with $x_i\in H$ and on the other coordinates, and hardwire them together with the bit $\bigoplus_{j\ne i}f(x_j)$}

\item \transfer\label{ex:pseudorandom:condenser} Lossless condensers. A function $\operatorname{Con}\colon\{0,1\}^n\times\{0,1\}^d\to\{0,1\}^m$, $m\ge k+d$, is a \emph{lossless condenser} for min-entropy $k$ with error $\eps$ if $\operatorname{Con}(X,Y)$ is $\eps$-close to some $(k+d)$-source for every $k$-source $X$ and independent uniform seed $Y$. Show that this holds if and only if its graph is a $({=}2^k,(1-\eps)2^d)$ vertex expander, and write the condition as a bound on lists; you may use that a $k$-source is a convex combination of flat ones when $2^k$ is an integer. (Ta-Shma, Umans, and Zuckerman; see \textcite[Proposition~6.33]{Vadhan2012}.)

\item \research\label{ex:pseudorandom:dp} Amplification in differential privacy. Several theorems of differential privacy state that a random operation applied to a private algorithm, such as running it on a random subsample or shuffling the reports of local randomizers, makes it more private. Can one of them be stated as a bound on $|\pseudorandomList_\Gamma(T,\theta)|$ for a natural $\Gamma$, tests, and threshold? If it can, then what do the translations of \cref{sec:pseudorandom:translations} give? If it cannot, then which feature of the framework fails? The question is open-ended.

\end{exercises}

\notesandrefs
The unified view is developed by \textcite{Vadhan2012}, whose chapters on expanders, codes, extractors, and generators each end with their list-decoding formulation. His Table~8.1 is the source of \cref{tab:pseudorandom:dictionary}, his chapter notes are the source of the attributions here that carry no other citation, and his unified nonconstructive object is the model for \cref{thm:pseudorandom:union}.

The dictionary's edges are marked by the randomness conductors of \textcite{CapalboReingoldVadhanWigderson2002} and the two-source extractors of \textcite{ChattopadhyayZuckerman2019}, both of which needed ideas from outside it.

\printchapterbib
\end{refsection}
